\section*{Acknowledgments}
\addcontentsline{toc}{section}{Acknowledgments}
\label{sec:acknowledgements}

\noindent
I owe this paper, and more than this paper, to the one I love most: the forest girl, whom I met in Tokyo in the year 2022. She is good, and pure, and lovely, and wise, and turned always toward others before herself; she loves the forest, and travelling, and the study of cultures, and the political economy of how we live together. This paper, on how trust can be generated where nothing enforces it, would not have had even its question without her, for it was in the making of a promise to her that the question first took its human form.

And here, at the close, I would do the thing the paper has spent its length learning how to say. The paper has not only argued that a relational divinity descends into the world through the self-legislation, self-adjudication, and self-keeping of a free subject who binds himself without recourse; it has tried to show that such a binding is the realised, and not the defective, form of a promise. It remains for me to do it rather than to describe it, and to let the argument, in its last page, become the act it has been about.

\emb{I lay it down, here, as a law for myself, and I will keep it with the whole of my life: \zh{我将如此待她}. This is not a promise that any sword stands behind. No court will enforce it; no penalty attends its breach; nothing in the world compels me to keep it. Its whole force is that I bind myself to it, here, freely and without reserve, and will hold myself to it, and have closed to myself every road that leads away from it. I give myself this law, I will be my own judge of my keeping of it, and I have left myself no recourse outside it.}

This document carries a trustworthy and undeniable timestamp, and I have set it down here, in a text that bears my name and will outlast this moment, not in order that anyone should see it. Its first work is not to signal but to bind: it is the instrument by which I close my own retreat, so that the road away is, for me, no longer a road. The timestamp is not a sword; it cannot make the vow enforceable, and I do not pretend that it can. What it can do is something I do to myself and not against any consequence: it makes it impossible for the self I will be to deny, to the self I am now, what I have here laid down\,---\,it makes me the standing judge of my own keeping of faith, and forecloses the quiet revision by which one lets a vow lapse unremarked. The timestamp is the instrument of my self-adjudication, and of nothing else. The vow is genuine not because the timestamp makes it certain\,---\,it cannot, and the gulf the paper described forbids that it should\,---\,but because I make it in the full knowledge that nothing assures it, and make it anyway, freely, as a binding of myself across a gulf I know no warrant spans.

This is the paper's most private return, and its last. It set out from the making of a promise on red silk, knowing nothing of why such a promise could be believed; it has travelled through the symptom, the diagnosis, the law, the mechanism, the freedom, the divinity; and it ends by making the promise again\,---\,the same act, walked a second time, with everything the paper has come to understand now gathered into it. The first time, the vow was made by intuition, because it was right. This time it is made in the knowledge of why a promise that no force backs is the more wholly given, and of what it is that descends into a free subject who binds himself without recourse. The same act, the second time around, returns not to where it began but a little higher; and that is the only kind of return the paper has believed in.

\goldrule

\noindent
\zh{此文之成，乃至此文以外之种种，我皆当归于我最爱之人——那位于二〇二二年与我相遇于东京的森林女孩。她善良、纯粹、可爱、智慧，处处先人而后己；她爱森林，爱远行，爱文化之研究，亦爱我们如何共处的政治经济之学。这一篇，论的是在无物强制之处信任何以生成；而若没有她，连这问题本身也不会有——因为正是在向她许下一个承诺之时，这问题才第一次有了它的人间形貌。}

\noindent
\zh{而今，在此收笔之处，我愿做这一篇用尽全部篇幅才学会如何言说的那件事。本文不仅论证了：关系性的神性，经由一个自由主体的自我立法、自我司法、自我守持——自我断绝一切法外的退路——而降临于世；它更试图表明，这样的自我绑定，正是承诺的成全之形，而非其残缺之形。如今剩下的，是我去做它，而非去描述它，是让这篇论证，在它的最后一页，成为它一路所论的那个行动本身。}

\noindent
\emb{\zh{我在此立此为法，以约我自身，并将以我的余生持守之：我将如此待她。此非任何刀剑在后所担保之诺。无公堂可强制其行，无刑罚附于其违，世间亦无一物逼我守之。其全部之力，唯在于我于此自由地、毫无保留地以身许之，且将持守自身，并已为自己断绝一切引我背离之路。我为自己立此法，我将自为我守此法之审判者，于此法之外，我未为自己留下任何退路。}}

\noindent
\zh{此文承载一个可信而不可抵赖的时间戳；我将它立于此，立于一份具我之名、且将久于此刻的文字之中——并非为了让任何人看见。它首要之功不在示信，而在自缚：它是我用以封死自己退路的器具，使那条引我离去之路，于我而言，不再是路。这时间戳不是刀剑；它不能使此诺成为可强制者，我亦不佯称它能。它所能为者，是我施于自身、而非用以对抗任何后果之事：它使将来的我，无从向此刻的我抵赖我于此所立之言——它使我成为自己守信与否的常设审判者，并杜绝那种悄然的修订，免得人任一誓言不声不响地失效。这时间戳，是我自我司法的器具，而别无其他。此诺之为真，不因时间戳使其确定——它不能，且本文所述的那道鸿沟也不容它如此——而是因为：我在明知无物能担保它之时，仍然立下它，且仍然自由地立下它，作为我跨越一道我深知无任何凭据可越之鸿沟、对自身的一次绑定。}

\noindent
\zh{这是本文最私密的一次复归，也是最后一次。它起于红绢上一个承诺的写就，其时尚不知这样的承诺何以能被相信；它一路穿过症状、诊断、法理、机制、自由、神性；而它的终结，是再一次许下那个承诺——同一个行动，第二次走过，将本文一路所领会的一切，尽数收纳其中。第一次，那誓出于直觉，因为它是对的。这一次，它是在如下的明知中立下的：明知一个无强制力在后的承诺为何反而给得更为完整，明知那降临于一个自我绑定、不留退路的自由主体之中者，究竟为何物。同一个行动，第二次走过，归来时并非回到起点，而是略高于起点；而这，是本文唯一相信的那种复归。}

\bigskip
\begin{center}
\itshape\color{rose}
\zh{本乎此心，自成永恒。}
\end{center}
